% Discussion v2 (20 Nov 2026). Every statement is backed by docs/results_ledger.md. Conclusion (about 70 words) still to write.

\section{Discussion and limitations}
\label{sec:discussion}
Sea surface temperature and sea level let a per-cell linear regression remove much of the climatological error; the U-Net ensemble adds a further $1.8$ to $4.9\%$ against GLORYS12, almost all at 50--200~m, and nothing useful near the surface or below 300~m. We read this as a measurement of what a deep network adds beyond a linear baseline under a protocol fixed before the test years, not as an accuracy record.

The reference matters more. The estimates are trained on GLORYS12, which is warmer than Argo in the thermocline, and they inherit the offset. The demeaned ensemble ranks last against GLORYS12 and first among the learned models against Argo, but one constant per depth, fitted on validation-window Argo, gives the frozen ensemble nearly the same Argo gain, so we read the demeaned advantage as bias reduction. The constants improve the frozen ensemble, the linear regression and GLORYS12 itself (by 5--6\%), but not the demeaned ensemble, which was already bias-reduced; the frozen ensemble keeps a smaller advantage of 3--5\% over the corrected linear regression. XGBoost, which receives the day of year, was better than the ensemble at 30--50~m in January--May of six of seven years (point estimates, no intervals); a season input is a hypothesis we did not test.

\emph{Limitations.} (i)~The target is a reanalysis that assimilates Argo, so the Argo comparison is not independent. (ii)~Intervals cover day-sampling variability only; Test~C is one season, 135 days (about 4.5 blocks) and 507 profiles. (iii)~The correction, its extension and the XGBoost analysis were specified after the test periods had been scored. (iv)~The single U-Net is a member of the ensemble, and the Andaman Sea advantage seen in validation did not replicate. (v)~Tests~B and C use interim GLORYS12; one domain, one reanalysis.
